Stuttering is a complex, multifactorial disorder characterized by atypical disruptions in the forward flow of speech~\cite{Smith_and_Weber}. It affects approximately 1 percent of the population and has a significant negative impact on all aspects of an individual’s life including social, educational, emotional, and vocational~\cite{craig_and_Tran,Koedoot_2011}.
% (Craig, Blumgart, and Tran, 2009; Koedoot, Bouwmans, Franken, and Stolk, 2011)
A stuttering assessment is conducted by a licensed speech-language pathologist. It includes measurement of the impact of stuttering on the person’s quality of life as well as analysis of the individual’s speech to determine stuttering severity. Unfortunately, over the past decades, speech-language pathologists have consistently reported limited competence in working with individuals who stutter 
% (e.g., 
% Brisk et al., 1997; Gabel, 2014; Kelly et al., 1997, St. Louis and Durrenberger, 1993)
(e.g.,~\cite{Brisk_1997,Gabel_2013,kelly_1997,stLouis_1994}
, with less than 5 percent of licensed professionals in the United States reporting expertise working with individuals who stutter~\cite{Coalson_2016,kelly_1997}.
% (Coalson et al., 2016; Kelly et al., 1997)
Another challenge is that the process of speech transcription and disfluency annotation is labor-intensive, as speech-language pathologists primarily transcribe and annotate speech disfluencies manually.  

A potential approach to address the above challenges is to use deep learning models to serve as an initial screening for stuttering, which could be deployed on edge devices such as smartphones or laptops. 
Early work~\cite{sheikhStutterNet,Kourkounakis2021FluentNetED}
has used simple deep learning models to detect stuttering behavior from speech utterances.
However, the main disadvantage is that the system is highly dependent on the scale of training data.
Unfortunately, the largest publicly available stuttering dataset is SEP-28K~\cite{lea_Sep28k}, which only contains 15.6 hours of utterance level labeled data, which is relatively small compared to standard datasets like LibriSpeech~\cite{librispeech}, which has a total 960 hours for training.

A promising recent approach in the speech community to deal with limited labeled data is Self-Supervised Learning (SSL).
Speech SSL models are first pretrained on a large amount of untranscribed speech and then finetuned on a smaller amount of transcribed/annotated speech for specific downstream tasks.
Due to the task-agnostic nature of the pretraining, speech SSL models demonstrate high generalizability across different speech processing tasks~\cite{yang21c_superb}.
Hence, speech SSL models are regarded as foundation models in many applications nowadays, such as Automatic Speech Recognition~\cite{baevski2021wav2vec-u,liu2022wav2vec-u2} and Speaker Verification~\cite{Peng_SSL_sv_2023}.
Due to the advantage of reducing the need for disfluency-labeled data, speech SSL models are a promising approach to tackling stuttered speech detection problems.
Specifically, prior works~\cite{Payal_iccasp_Siamese,bayerl22b_interspeech,bayerl23_interspeech} utilize Wav2vec 2.0~\cite{wav2vec2} as their model backbone for stutterred speech detection.

% mention that previous work do not focus on word level stutter evaluation
Despite progress in this field, previous work has mainly focused on utterance-level stuttering detection.
However, this is too coarse for clinical application as stuttering-like disfluencies are also present in typically fluent individuals who are not diagnosed as a Person who Stutters~(PWS).
Clinically, stuttering is diagnosed based on meeting a certain frequency threshold of stuttering-like disfluencies.  
Hence, we are interested in developing a more fine-grained stuttering detection model.
To our best knowledge, ~\cite{harvill22_interspeech} is the only previous work that investigates the task of frame-level stuttering detection.
However, they evaluated their model by removing the segments that are predicted as stuttering and conducted a user study on Amazon Mechanical Turk, which was not assessed on clinical data collected and annotated by speech-language pathologists (SLPs).

In our work, we first collect a set of speech recordings with stuttering events annotated by speech-language pathologists and use it as a word-level stuttering detection evaluation benchmark.
Then we propose a word-level stuttering detection model that utilizes WavLM~\cite{Chen2021WavLMLS} as our backbone.
Following~\cite{harvill22_interspeech}, we first pretrain on LibriSpeech with synthetic disfluency augmentations and then finetune our model on SEP-28K dataset.
We show that our model not only shows strong performance on utterance-level stuttering detection on SEP28K but also outperforms prior work on word-level stuttered speech benchmark by a large margin.
In addition, we study the effect of the utilization of Speech SSL models on word-level stuttering detection with extensive experiments.
% In addition to tackle data deficiency problem by using Speech SSL model, we propose to use voice conversion models to scale up the SEP28K as we find out that voice conversion model can preserve the characteristic of stuttering while changing speaker characteristics.
% Our model further improve on word level stuttering detection.
% Last but not least, we conduct extensive ablation studies on the performance of using Speech SSL models for word level stuttering detection.
In conclusion, our contribution can be summarized as the following:
\begin{enumerate}
    \item We are the first to focus on word-level stuttering detection, which is more closely aligned to clinical applications.
    \item We propose a word-level stuttering detection model that achieves state-of-the-art performance.
    \item We study the effect of utilizing Speech Self-supervised models on word-level stuttering detection with clinical data evaluation and extensive ablations.
    % \item We open a new direction for incorporating voice conversion in stuttered speech detection research to tackle the problem of label scarcity which could also be regarded as as way to preserve stuttering patients' privacy.
\end{enumerate}
We open source our code on Github\footnote{\url{www.github.com/atosystem/SpeechSSLStutterDetect}}.
